\section{Context, effort and persistent instructions}
% Block 3 of PLAN.md, 20 min.  Message: the context is the model's only working
% memory, it is finite and degrades as it fills; what must survive goes to disk.
% This is the block to shorten if time runs out.

% ------------------------------------------------------------------------------
\begin{frame}{The context window}{The model's only working memory}
  \vspace{2ex}
  \begin{center}
    \begin{tikzpicture}[seg/.style={draw=unibasWhite,line width=0.8pt,minimum height=8mm,
        inner sep=0pt,font=\tiny,text=unibasWhite,align=center}]
      \def\L{11}
      \draw[unibasLightGrey,line width=0.6pt] (0,-0.4) rectangle (\L,0.4);
      \node[lbl,anchor=south east] at (\L,0.45) {limit: $\sim 10^5$--$10^6$ tokens};
      \node[seg,fill=unibasAnthracite,minimum width=8mm,anchor=west] (a) at (0,0) {system};
      \node[seg,fill=unibasPetrol,minimum width=12mm,right=0pt of a] (b) {instruction\\files};
      \onslide<2->{\node[seg,fill=figHuman!70,minimum width=10mm,right=0pt of b] (c) {your\\messages};}
      \onslide<3->{\node[seg,fill=figTool,minimum width=24mm,right=0pt of c] (d) {files read};
                   \node[seg,fill=figTool!70,minimum width=26mm,right=0pt of d] (e) {tool output};}
      \onslide<4->{\node[seg,fill=figModel!70,minimum width=16mm,right=0pt of e] (f) {model's\\answers};}
    \end{tikzpicture}
  \end{center}
  \vspace{1ex}

  \begin{itemize}
    \item Everything the model knows during this session, and nothing else
    \item<3-> Every file read and every command output adds to it
    \item<4-> The limit is a few books; a long agent session reaches it
  \end{itemize}
\end{frame}

% ------------------------------------------------------------------------------
\begin{frame}{More context is not better context}{Context rot}
  \begin{itemize}
    \item Accuracy drops as the context grows, long before the limit
    \item Irrelevant text distracts; a wrong turn stays in memory
    \item<2-> Near the limit, the harness summarises the history: lossy
  \end{itemize}

  \vspace{3ex}
  \onslide<3->{\centering\ckeybox[0.6\linewidth]{One task, one fresh session}\par}

  \slidecite{\citeContextRot}
\end{frame}

% ------------------------------------------------------------------------------
\begin{frame}[fragile]{Context hygiene}{Point, do not paste}
\only<1>{%
  \vspace{0.5ex}{\scriptsize\color{unibasGrey}Instead of this \dots}}
\begin{prompt}
It doesn't work, here is the output: [3000 lines of solver log]
\end{prompt}
\begin{onlyenv}<2->
\vspace{-0.5ex}{\scriptsize\color{unibasGrey}\dots\ this:}
\begin{prompt}
The Ramsey test in test/runtests.jl fails. Run only that testset,
read the failing assertion, and explain the cause before changing code.
\end{prompt}
\end{onlyenv}
  \begin{itemize}
    \item<3-> Name files and functions; the agent reads what it needs
    \item<3-> After a wrong turn, start a new session rather than arguing
  \end{itemize}
\end{frame}

% ------------------------------------------------------------------------------
\begin{frame}{Thinking costs tokens}{Effort and token economics}
  \begin{equation*}
    \text{cost} \approx N_{\mathrm{in}}\,c_{\mathrm{in}} + N_{\mathrm{out}}\,c_{\mathrm{out}}
  \end{equation*}

  \begin{itemize}
    \item $N_{\mathrm{in}}$: tokens read; every agent step re-reads the context
    \item $N_{\mathrm{out}}$: tokens written, including hidden ``thinking''; $c$: price per token
    \item<2-> \emph{Effort} setting: how much the model thinks before answering
    \item<3-> Low for mechanical edits, high for debugging and derivations
  \end{itemize}

  \vspace{1ex}
  \onslide<4->{{\small Scale: a research-level physics problem for under 3 USD}}

  \slidecite{\citeTheoPhysEval}
\end{frame}

% ------------------------------------------------------------------------------
\begin{frame}[fragile]{Write it down: instruction files}{The agent forgets; the file does not}
\begin{terminal}[title={\hfill\tiny\color{unibasGrey}\smallcapsline{DrivenSpin/AGENTS.md (excerpt)}}]
## Physics conventions
- ħ = 1; times in μs, angular frequencies in rad/μs
- Basis (|e⟩, |g⟩); H = (Δ/2)σz + (Ω/2)σx with Δ = ω_q - ω_d
- Never change a convention without asking

## How to work
- Run the tests: julia --project -e 'using Pkg; Pkg.test()'
- Every new function gets a test with a physical check
- Never weaken or delete a test to make it pass: ask first
- Never write to data/exp_raw
\end{terminal}
  \begin{itemize}
    \item<2-> Read automatically at the start of every session
    \item<3-> \cmd{AGENTS.md}, \cmd{CLAUDE.md}, \dots: the file name depends on the tool,
          the idea does not
  \end{itemize}
\end{frame}

% ------------------------------------------------------------------------------
\begin{frame}[fragile]{Skills: procedures on demand}{Written once, loaded when relevant}
\begin{terminal}[title={\hfill\tiny\color{unibasGrey}\smallcapsline{skills/fit-ramsey/SKILL.md}}]
---
name: fit-ramsey
description: Fit Ramsey fringes and report T2* with errors. Use when
  analysing Ramsey or free-induction-decay data.
---
1. Load the data with load_ramsey(path); never modify the raw file.
2. Fit P(t) = A exp(-t/T2*) cos(Δt + φ) + B with LsqFit.
3. Check: residuals show no structure; T2* ≤ 2 T1 if T1 is known.
4. Save the fit with produce_or_load; plot data and fit; report T2* ± σ.
\end{terminal}
  \begin{itemize}
    \item<2-> Only \cmd{name} and \cmd{description} sit in the context: cheap
    \item<3-> The body is loaded when a task matches: your lab's procedures, on demand
  \end{itemize}
\end{frame}

% ------------------------------------------------------------------------------
\begin{frame}[fragile]{A good prompt is a specification}
\begin{onlyenv}<1>
\begin{prompt}
Simulate a qubit with decoherence.
\end{prompt}
\end{onlyenv}
\begin{onlyenv}<2->
\begin{prompt}
Add T1 and T2 decay to the driven-spin model in src/DrivenSpin.jl, using
the Lindblad equation with the conventions in AGENTS.md.
Before writing code, propose a plan and the tests: trace preservation,
positivity, and free decay exp(-t/T1) of the excited population.
Do not change the existing Hamiltonian or its tests.
\end{prompt}
\end{onlyenv}
  \begin{itemize}
    \item<2-> \textbf{Goal} and \textbf{context}: which file, which conventions
    \item<3-> \textbf{Verification}: the tests that define ``done''
    \item<4-> \textbf{Constraints}: what must not change; plan before code
  \end{itemize}
\end{frame}

% ------------------------------------------------------------------------------
\begin{frame}{A workflow that works}
  \vspace{3ex}
  \begin{center}
    \begin{tikzpicture}[node distance=5mm]
      \node[st=figModel] (e) {\textbf{Explore}\\{\scriptsize read the code}};
      \node[st=figModel,right=of e] (p) {\textbf{Plan}\\{\scriptsize you approve}};
      \node[st=figModel,right=of p] (i) {\textbf{Implement}\\{\scriptsize small steps}};
      \node[st=figTool,right=of i] (v) {\textbf{Verify}\\{\scriptsize tests, plots}};
      \node[st=figHuman,right=of v] (c) {\textbf{Commit}\\{\scriptsize you read the diff}};
      \foreach \a/\b in {e/p,p/i,i/v,v/c} \draw[flow] (\a) -- (\b);
    \end{tikzpicture}
  \end{center}
  \vspace{2ex}

  \begin{itemize}
    \item<2-> Most failures come from skipping \emph{plan} or \emph{verify}
    \item<3-> \emph{Commit} needs version control: next section
  \end{itemize}
\end{frame}
